\documentclass[a4paper, serif, 10pt]{moderncv}

%% moderncv
% style and color
\moderncvstyle{banking}
\moderncvcolor{black}

% renew moderncv command to adapt for CJK colon
\renewcommand*{\cvitem}[3][.25em]{%
  \ifstrempty{#2}{}{\hintstyle{#2}：}{#3}%
  \par\addvspace{#1}}

\renewcommand*{\cvitemwithcomment}[4][.25em]{%
  \savebox{\cvitemwithcommentmainbox}{\ifstrempty{#2}{}{\hintstyle{#2}：}#3}%
  \setlength{\cvitemwithcommentmainlength}{\widthof{\usebox{\cvitemwithcommentmainbox}}}%
  \setlength{\cvitemwithcommentcommentlength}{\maincolumnwidth-\separatorcolumnwidth-\cvitemwithcommentmainlength}%
  \begin{minipage}[t]{\cvitemwithcommentmainlength}\usebox{\cvitemwithcommentmainbox}\end{minipage}%
  \hfill% fill of \separatorcolumnwidth
  \begin{minipage}[t]{\cvitemwithcommentcommentlength}\raggedleft\small\itshape#4\end{minipage}%
  \par\addvspace{#1}}

%% page layout/margins
\usepackage[top=2.5cm, bottom=2.5cm, left=1.5cm, right=1.5cm]{geometry}

\nopagenumbers{}

%% Babel config for English language
\usepackage[english]{babel}

%% fontspec
\usepackage{fontspec}

\IfFontExistsTF{Linux Libertine}{
  \setmainfont[Ligatures={TeX, Common}, Numbers=Lining, ItalicFont=Linux Libertine]{Linux Libertine}
}{}
\IfFontExistsTF{Linux Libertine O}{
  \setmainfont[Ligatures={TeX, Common}, Numbers=Lining, ItalicFont=Linux Libertine O]{Linux Libertine O}
}{}

%% CTeX
% CJK support, used to show CJK characters in the resume
%
% - fontset=none: disable builtin fontset but instead set the CJK font manually
% - heading=false: disable ctex heading
% - punct=kaiming: use kaiming punctuations styles for CJK
% - scheme=plain: use plain scheme, do not override \normalsize font size
% - space=auto: space settings for CJK characters
%
% ref:
% - http://ctan.mirrorcatalogs.com/language/chinese/ctex/ctex.pdf
\usepackage[UTF8, heading=false, punct=kaiming, scheme=plain, space=auto]{ctex}

\IfFontExistsTF{Noto Serif CJK SC}{
  \setCJKmainfont{Noto Serif CJK SC}
}{}
\IfFontExistsTF{Noto Sans CJK SC}{
  \setCJKsansfont{Noto Sans CJK SC}
}{}

%% line spacing
\usepackage{setspace}
\setstretch{1.125}

%% URL styling - use normal text instead of monospace
\urlstyle{same}

% Auto-underline all links
% Defer until \begin{document} so hyperref is already loaded
\AtBeginDocument{%
  \let\oldhref\href
  \renewcommand{\href}[2]{\oldhref{#1}{\underline{#2}}}%
}

%% Patch \httplink and \httpslink to handle full URLs
% The original moderncv macros always prepend http:// or https:// to URLs.
% This causes issues when the URL already has a protocol (e.g., https://example.com)
% resulting in malformed URLs like https://https://example.com.
% This patch checks if the URL already contains "://" and uses it directly if so.
% Note: We use \str_set:Nx (x-expansion) to expand macros like \@homepage first.
\ExplSyntaxOn
\str_new:N \g_yamlresume_url_str

\RenewDocumentCommand{\httpslink}{O{}m}{
  \str_set:Nx \g_yamlresume_url_str {#2}
  \str_if_in:NnTF \g_yamlresume_url_str {://}
    { \url{#2} }
    { \url{https://#2} }
}

\RenewDocumentCommand{\httplink}{O{}m}{
  \str_set:Nx \g_yamlresume_url_str {#2}
  \str_if_in:NnTF \g_yamlresume_url_str {://}
    { \url{#2} }
    { \url{http://#2} }
}
\ExplSyntaxOff

\name{陳雅婷}{}
\title{資深產品經理 | 金融科技與數碼解決方案}
\phone[mobile]{+852 9123 4567}
\email{ngating.chan@example.com}
\homepage{https://www.linkedin.com/in/ngatingchan}

\address{中國香港，香港島，香港，香港鰂魚涌英皇道 979 號}{}{}

\extrainfo{{\small \faLinkedin}\ \href{https://www.linkedin.com/in/ngatingchan}{@ngatingchan} {} {} {} • {} {} {}
{\small \faMedium}\ \href{https://medium.com/@ngatingchan}{@ngatingchan}}

\begin{document}

\maketitle

\section{簡介}

\cvline{}{擁有超過 8 年產品管理經驗的資深產品經理，專注於金融科技、電子商貿及流動應用程式領域。擅長由 0 到 1 建立產品，並通過數據驅動決策、用戶研究及跨團隊協作，持續提升產品價值及用戶體驗。
%
\begin{itemize}
\item 主導多個大型數碼產品項目，涵蓋流動銀行、數碼支付及零售電商，累計服務用戶超過 200 萬
\item 熟悉敏捷開發、精實啟動及 OKR 管理，能夠有效協調工程、設計、市場及業務團隊
\item 擁有數據分析及 A/B 測試實戰經驗，利用 SQL、Amplitude、Mixpanel 等工具制定產品策略
\item 具備國際視野及本地市場洞察，能靈活應對瞬息萬變的業務需求
\end{itemize}}
\section{教育背景}

\cventry{2010年9月–2014年7月}
        {學士，工商管理學士（資訊系統），成績：3.7}
        {香港中文大學}
        {\url{https://www.cuhk.edu.hk/}}
        {}
        {\begin{itemize}
\item 主修資訊系統，副修市場學，畢業論文探討流動支付在香港的普及因素
\item 擔任學生會資訊科技幹事，統籌校園數碼轉型項目
\item 獲院長嘉許名單（Dean's List）兩次
\end{itemize}
\textbf{課程}：產品管理與開發、數據庫管理、電子商貿策略、市場研究與分析、資訊系統審計、企業架構}
\section{工作經歷}

\cventry{2022年4月至今}
        {高級產品經理}
        {智融香港有限公司}
        {\url{https://www.smartfintech.com.hk}}
        {}
        {\begin{itemize}
\item 主導新一代流動銀行應用程式的產品規劃，由研究、設計、開發至上市全程管理，用戶評分提升至 4.8 星
\item 與工程、設計、風控及合規團隊緊密合作，確保產品符合金管局監管要求及銀行安全標準
\item 制定並執行產品路線圖，利用 OKR 框架與持份者對齊目標，按季度交付 15 項以上功能更新
\item 建立數據分析機制，透過 A/B 測試優化入門流程，令新用戶轉換率提升 32\%
\item 帶領 3 名產品助理及 2 名產品設計師，建立以用戶為中心的產品文化
\end{itemize}
\textbf{關鍵字}：金融科技、流動銀行、產品路線圖、A/B 測試、敏捷開發、持份者管理}

\cventry{2018年7月–2022年3月}
        {產品經理}
        {數碼支付香港有限公司}
        {\url{https://www.digitalpay.com.hk}}
        {}
        {\begin{itemize}
\item 負責數碼錢包及電子支付產品的端到端交付，包括 QR Code 支付、P2P 轉賬及商戶管理系統
\item 與超過 50 間合作商戶進行需求訪談，設計並推出商家促銷及會員積分功能，帶動交易量增長 45\%
\item 利用 SQL 及 Amplitude 分析用戶行為，提出並落實減低支付失敗率的措施，成功將失敗率由 3.5\% 降至 1.8\%
\item 參與 Scrum 團隊，擔任 Product Owner，撰寫用戶故事及驗收準則，確保每個 Sprint 按時交付
\end{itemize}
\textbf{關鍵字}：數碼支付、Fintech、用戶研究、SQL、Product Owner、電子錢包}

\cventry{2015年8月–2018年6月}
        {助理產品經理}
        {環球電商方案有限公司}
        {\url{https://www.globalecom.io}}
        {}
        {\begin{itemize}
\item 協助高級產品經理管理 B2B 電子商貿平台的功能開發，涵蓋商品管理、訂單處理及物流追蹤模組
\item 撰寫產品需求文件（PRD），與工程師及設計師協作進行功能驗收，確保產品質素
\item 定期分析網站流量及銷售數據，為團隊提供營運建議，推動平均訂單價值提升 18\%
\item 協調客戶服務部門處理用戶反饋，整理常見問題並優化自助服務流程
\end{itemize}
\textbf{關鍵字}：B2B、電子商貿、PRD、數據分析、客戶成功}
\section{語言}

\cvline{粵語}{母語或雙語水平 \hfill \textbf{關鍵字}：母語}
\cvline{普通話}{母語或雙語水平 \hfill \textbf{關鍵字}：母語}
\cvline{英語}{完全專業水平 \hfill \textbf{關鍵字}：IELTS Overall 8.0、TOEIC 950}
\section{專業技能}

\cvline{產品管理}{專家 \hfill \textbf{關鍵字}：產品策略、路線圖規劃、用戶故事、敏捷／Scrum、持份者管理、競品分析}
\cvline{數據分析}{高級 \hfill \textbf{關鍵字}：SQL、Amplitude、Mixpanel、Google Analytics、A/B Testing}
\cvline{用戶體驗設計}{中級 \hfill \textbf{關鍵字}：Figma、用戶訪談、可用性測試、Wireframe、用戶流程圖}
\cvline{商業工具}{專家 \hfill \textbf{關鍵字}：Jira、Confluence、Notion、Tableau、Excel}
\section{榮譽}

\cventry{2021年10月}
        {香港金融科技協會}
        {香港金融科技發展獎（最佳流動支付體驗）}
        {}
        {}
        {憑藉數碼錢包產品的創新功能和用戶體驗設計，獲頒「最佳流動支付體驗」獎項，表揚團隊在支付流程簡化及安全合規方面的卓越表現。}
\section{證書}

\cventry{2019年6月}
        {Scrum Alliance}
        {Certified Scrum Product Owner (CSPO)}
        {\url{https://www.scrumalliance.org/}}
        {}
        {}

\cventry{2020年8月}
        {香港產品管理學會}
        {高級產品管理專業證書}
        {\url{https://www.hkpm.org.hk}}
        {}
        {}
\section{出版刊物}

\cventry{2022年3月}
        {數碼產品經理必須掌握的數據分析方法}
        {Medium}
        {\url{https://medium.com/@ngatingchan/data-driven-product-management}}
        {}
        {分享如何利用 SQL 及 A/B 測試制定產品決策，並以實際案例說明數據分析在金融科技產品中的應用。}
\section{項目}

\cventry{2023年1月–2023年7月}
        {一款專為年輕用戶設計的互動理財學習應用程式，利用遊戲化元素提升金融知識普及率。}
        {投資知識學習平台「理財新手村」}
        {\url{https://www.finlearn.hk}}
        {}
        {\begin{itemize}
\item 針對 18-30 歲用戶缺乏理財知識的痛點，主導產品概念開發及用戶驗證
\item 與教育機構及設計顧問合作，設計 12 個互動課程及投資情境模擬
\item 上線三個月內錄得超過 10,000 名註冊用戶，課程完成率達 68\%
\end{itemize}
\textbf{關鍵字}：金融教育、遊戲化、用戶研究、產品設計}
\section{興趣愛好}

\cvline{閱讀}{商業書籍、行為經濟學、科技趨勢}
\cvline{戶外活動}{行山、露營、跑步}
\section{志願服務}

\cventry{2020年1月–2023年12月}
        {義務導師}
        {香港產品經理協會}
        {\url{https://www.hkpmassociation.org}}
        {}
        {\begin{itemize}
\item 為轉職產品管理人士提供一對一職業指導，分享求職及面試技巧
\item 定期主持工作坊，主題包括「從 0 到 1 建立產品」及「數據驅動決策」
\item 協助協會舉辦年度研討會，邀請業界領袖分享最新產品管理趨勢
\end{itemize}}

\end{document}